% ============================================================
% CHAPITRE 13 — Création de supports publicitaires avec Photoshop
% ============================================================

\chapter{Création de supports publicitaires avec Photoshop}

\noindent
\textit{Nous voici au dernier chapitre de la deuxième partie, consacrée à la maîtrise d'Adobe Photoshop. Au fil des chapitres précédents, nous avons appris à créer des documents, à organiser des calques, à sélectionner, à détourer, à retoucher des photographies, à réaliser des photomontages, à intégrer du texte et à appliquer des effets graphiques. Toutes ces compétences techniques sont désormais acquises. Il reste à les mettre en œuvre dans des projets concrets : la création de supports publicitaires.}

\vspace{1em}

\noindent
La \terme{création de supports publicitaires} est l'une des applications les plus courantes de Photoshop en agence, en entreprise, en freelance. Affiches, flyers, visuels pour les réseaux sociaux, bannières web : tous ces supports sont conçus avec Photoshop (parfois complétés par Illustrator ou InDesign). Ils servent à promouvoir un produit, un service, un événement, une marque. Ils doivent être efficaces : attirer l'attention, transmettre un message, susciter une action.

\vspace{1em}

\noindent
Mais un support publicitaire n'est pas seulement une composition esthétique : c'est un \terme{outil de communication} qui doit atteindre des objectifs précis. Il doit respecter les contraintes du support (format, résolution, colorimétrie), il doit s'adapter au public (âge, milieu, habitudes), il doit respecter la charte graphique du client, il doit être décliné sur plusieurs formats, il doit être exporté dans les bons formats. C'est cette dimension professionnelle que nous allons explorer dans ce chapitre.

\vspace{1em}

\noindent
Ce chapitre a pour objectif de te faire maîtriser la création de supports publicitaires avec Photoshop. Nous commencerons par une méthodologie de création d'une affiche. Nous verrons ensuite comment concevoir un flyer commercial, comment créer des visuels pour les réseaux sociaux, comment réaliser une affiche événementielle, comment créer une bannière publicitaire. Nous apprendrons à intégrer des photographies et des textes, à organiser la hiérarchie visuelle, à décliner un visuel en plusieurs formats, à contrôler la qualité, et à exporter les fichiers finaux. À la fin de ce chapitre, tu sauras créer des supports publicitaires professionnels.

% ============================================================
\section{Méthodologie de création d'une affiche}
% ============================================================

\subsection{Pourquoi une méthodologie ?}

La création d'une affiche ne s'improvise pas. Elle repose sur une \terme{méthodologie} qui permet de :

\begin{itemize}
    \item \terme{Structurer} la réflexion.
    \item \terme{Gagner} du temps.
    \item \terme{Éviter} les erreurs.
    \item \terme{Produire} un résultat de qualité.
    \item \terme{Communiquer} avec le client.
    \item \terme{Reproduire} le processus sur d'autres projets.
\end{itemize}

La méthodologie est la même que celle du processus de création graphique (voir chapitre 1), mais adaptée au contexte spécifique de l'affiche.

\subsection{Les étapes de création d'une affiche}

\subsubsection{Étape 1 : Le brief}

Comprendre le besoin du client :

\begin{itemize}
    \item \terme{Contexte} : qui est le client, que fait-il, quel est son marché ?
    \item \terme{Objectif} : que veut-il obtenir ? (notoriété, ventes, information)
    \item \terme{Public} : à qui s'adresse l'affiche ? (âge, milieu, habitudes)
    \item \terme{Message} : quel est le message principal ?
    \item \terme{Contraintes} : format, budget, délai, supports, charte graphique.
    \item \terme{Références} : ce que le client aime, ce qu'il n'aime pas.
\end{itemize}

\subsubsection{Étape 2 : La recherche}

Rechercher des références et s'inspirer :

\begin{itemize}
    \item \terme{Analyse} du marché et de la concurrence.
    \item \terme{Recherche} de références visuelles (Behance, Pinterest, magazines).
    \item \terme{Création} d'un moodboard.
    \item \terme{Définition} de l'ambiance (couleurs, typographies, style).
\end{itemize}

\subsubsection{Étape 3 : La conception}

Concevoir l'affiche :

\begin{itemize}
    \item \terme{Définir} le concept.
    \item \terme{Esquisser} plusieurs pistes (croquis papier).
    \item \terme{Sélectionner} la piste la plus pertinente.
    \item \terme{Définir} la hiérarchie visuelle.
\end{itemize}

\subsubsection{Étape 4 : La création}

Créer l'affiche dans Photoshop :

\begin{itemize}
    \item \terme{Créer} le document aux bonnes dimensions (A3, A4, etc.).
    \item \terme{Importer} les images, les détourer si nécessaire.
    \item \terme{Composer} : placer les éléments selon la hiérarchie.
    \item \terme{Appliquer} les couleurs, les typographies, les effets.
    \item \terme{Ajuster} les détails.
\end{itemize}

\subsubsection{Étape 5 : La présentation}

Présenter l'affiche au client :

\begin{itemize}
    \item \terme{Expliquer} la démarche, les choix, les bénéfices.
    \item \terme{Présenter} sur écran et/ou sur papier.
    \item \terme{Recueillir} les retours.
    \item \terme{Négocier} les ajustements.
\end{itemize}

\subsubsection{Étape 6 : La finalisation}

Finaliser l'affiche :

\begin{itemize}
    \item \terme{Intégrer} les modifications.
    \item \terme{Vérifier} la qualité (alignements, couleurs, résolution).
    \item \terme{Préparer} les fichiers pour la production.
    \item \terme{Exporter} dans le bon format.
\end{itemize}

\subsubsection{Étape 7 : La production}

Produire et diffuser l'affiche :

\begin{itemize}
    \item \terme{Envoyer} à l'imprimeur (si impression).
    \item \terme{Contrôler} le résultat final.
    \item \terme{Diffuser} (affichage, réseaux sociaux, web).
\end{itemize}

\subsubsection{Étape 8 : L'évaluation}

Évaluer le résultat :

\begin{itemize}
    \item \terme{Vérifier} que les objectifs sont atteints.
    \item \terme{Recueillir} les retours du client, du public.
    \item \terme{Archiver} les fichiers.
    \item \terme{Tirer} les leçons du projet.
\end{itemize}

\subsection{Les contraintes de l'affiche}

\subsubsection{Le format}

Le format de l'affiche dépend de son usage :

\begin{itemize}
    \item \terme{A3} (29,7 × 42 cm) : affichage intérieur, vitrines.
    \item \terme{A2} (42 × 59,4 cm) : affichage intérieur, commerces.
    \item \terme{A1} (59,4 × 84,1 cm) : affichage extérieur, petits panneaux.
    \item \terme{A0} (84,1 × 118,9 cm) : affichage extérieur, grands panneaux.
    \item \terme{4 × 3 m} : affichage grand format (panneaux publicitaires).
\end{itemize}

\subsubsection{La résolution}

La résolution dépend de la distance de lecture :

\begin{itemize}
    \item \terme{300 ppi} : pour les affiches vues de près (A3, A4).
    \item \terme{150 ppi} : pour les affiches vues de loin (A2, A1).
    \item \terme{72 ppi} : pour les affiches grand format (4 × 3 m).
\end{itemize}

\subsubsection{La colorimétrie}

Pour l'impression, travailler en \terme{CMJN}. Pour l'écran, travailler en \terme{RVB}.

\subsubsection{Les marges et les fonds perdus}

Prévoir :

\begin{itemize}
    \item \terme{Marges de sécurité} : 5 mm des bords.
    \item \terme{Fonds perdus} : 3 à 5 mm au-delà des bords.
    \item \terme{Traits de coupe} : repères pour l'imprimeur.
\end{itemize}

\subsection{Les bonnes pratiques}

\begin{itemize}
    \item \terme{Commencer} par le brief.
    \item \terme{Faire} de la recherche.
    \item \terme{Esquisser} plusieurs pistes.
    \item \terme{Définir} la hiérarchie visuelle.
    \item \terme{Respecter} les contraintes techniques.
    \item \terme{Présenter} et argumenter.
    \item \terme{Vérifier} la qualité avant livraison.
    \item \terme{Archiver} les fichiers.
\end{itemize}

\subsection*{Conclusion de la section}

La \terme{méthodologie de création d'une affiche} suit huit étapes : brief, recherche, conception, création, présentation, finalisation, production, évaluation. Elle permet de structurer la réflexion, de gagner du temps, d'éviter les erreurs. Les contraintes de l'affiche sont le format, la résolution, la colorimétrie, les marges et fonds perdus. Il est essentiel de commencer par le brief, de faire de la recherche, d'esquisser, de définir la hiérarchie, de respecter les contraintes, de vérifier la qualité, d'archiver.

Maintenant que nous connaissons la méthodologie, nous allons, dans la section suivante, apprendre à \terme{concevoir un flyer commercial}.

% ============================================================
\section{Conception d'un flyer commercial}
% ============================================================

\subsection{Qu'est-ce qu'un flyer ?}

\begin{quote}
\textbf{Un flyer (ou tract) est un support de communication imprimé, de petit format (A6, A5, A4), destiné à promouvoir un produit, un service, un événement, une marque. Il est distribué à la main, par courrier, ou déposé dans des lieux publics.}
\end{quote}

Le flyer est un support \terme{économique}, \terme{rapide}, \terme{efficace}. Il est utilisé par les commerces, les associations, les entreprises, les événements. Il doit attirer l'attention, transmettre un message, inciter à l'action.

\subsection{Les caractéristiques du flyer}

\subsubsection{Le format}

\begin{itemize}
    \item \terme{A6} (10,5 × 14,8 cm) : petit, économique, pour une information brève.
    \item \terme{A5} (14,8 × 21 cm) : courant, pour une information moyenne.
    \item \terme{A4} (21 × 29,7 cm) : grand, pour une information détaillée.
    \item \terme{DL} (10 × 21 cm) : format allongé, pour une information listée.
    \item \terme{Carré} (14 × 14 cm) : moderne, pour une information visuelle.
\end{itemize}

\subsubsection{Le recto-verso}

Un flyer peut être :

\begin{itemize}
    \item \terme{Recto seul} : une face, information brève.
    \item \terme{Recto-verso} : deux faces, information détaillée.
    \item \terme{Plié} : plusieurs volets, information structurée.
\end{itemize}

\subsubsection{Le papier}

\begin{itemize}
    \item \terme{Papier standard} : 135 g à 170 g, économique.
    \item \terme{Papier épais} : 250 g à 350 g, qualitatif.
    \item \terme{Papier recyclé} : écologique.
    \item \terme{Papier brillant} : moderne, pour les photos.
    \item \terme{Papier mat} : élégant, pour les textes.
\end{itemize}

\subsection{La structure d'un flyer}

\subsubsection{Le recto}

Le \terme{recto} est la face la plus visible. Il doit :

\begin{itemize}
    \item \terme{Attirer} l'attention (titre, visuel fort).
    \item \terme{Transmettre} le message principal.
    \item \terme{Donner} envie de lire le verso.
\end{itemize}

Éléments du recto :

\begin{itemize}
    \item \terme{Titre} : accroche, message principal.
    \item \terme{Visuel} : image forte, attractive.
    \item \terme{Logo} : identification de la marque.
    \item \terme{Accroche} : phrase courte, percutante.
\end{itemize}

\subsubsection{Le verso}

Le \terme{verso} donne les informations détaillées. Il doit :

\begin{itemize}
    \item \terme{Informer} : description, avantages, offres.
    \item \terme{Convaincre} : arguments, témoignages.
    \item \terme{Inciter} à l'action : appel à l'action, coordonnées.
\end{itemize}

Éléments du verso :

\begin{itemize}
    \item \terme{Description} : présentation du produit ou service.
    \item \terme{Avantages} : bénéfices pour le client.
    \item \terme{Offre} : prix, promotion, conditions.
    \item \terme{Coordonnées} : adresse, téléphone, site web, réseaux sociaux.
    \item \terme{Appel à l'action} : ce que le client doit faire.
    \item \terme{Mentions légales} : obligations légales.
\end{itemize}

\subsection{La conception d'un flyer commercial}

\subsubsection{Étape 1 : Définir le message}

Quel est le message principal ? Que doit retenir le lecteur ?

\subsubsection{Étape 2 : Définir le public}

À qui s'adresse le flyer ? Quel est son âge, son milieu, ses habitudes ?

\subsubsection{Étape 3 : Choisir le format}

Quel format est adapté au message et au public ? (A5, A6, DL, etc.)

\subsubsection{Étape 4 : Concevoir le recto}

Créer un recto attractif :

\begin{itemize}
    \item Titre fort.
    \item Visuel attractif.
    \item Logo.
    \item Accroche.
\end{itemize}

\subsubsection{Étape 5 : Concevoir le verso}

Créer un verso informatif :

\begin{itemize}
    \item Description.
    \item Avantages.
    \item Offre.
    \item Coordonnées.
    \item Appel à l'action.
    \item Mentions légales.
\end{itemize}

\subsubsection{Étape 6 : Vérifier la cohérence}

Le recto et le verso doivent être cohérents : mêmes couleurs, mêmes typographies, même style.

\subsection{Les bonnes pratiques}

\begin{itemize}
    \item \terme{Un seul message} principal.
    \item \terme{Hiérarchie} claire.
    \item \terme{Typographie} limitée (1 à 2 polices).
    \item \terme{Couleurs} cohérentes (2 à 4).
    \item \terme{Visuel} de qualité.
    \item \terme{Coordonnées} visibles.
    \item \terme{Appel à l'action} clair.
    \item \terme{Mentions légales} présentes.
    \item \terme{Vérifier} la lisibilité.
    \item \terme{Tester} sur papier avant impression.
\end{itemize}

\subsection*{Conclusion de la section}

Le \terme{flyer} est un support de communication imprimé, de petit format, destiné à promouvoir un produit, un service, un événement, une marque. Il peut être recto, recto-verso ou plié. Le recto attire l'attention, le verso informe. La conception suit plusieurs étapes : définir le message, définir le public, choisir le format, concevoir le recto, concevoir le verso, vérifier la cohérence. Il est essentiel d'avoir un seul message, une hiérarchie claire, une typographie limitée, des couleurs cohérentes, des coordonnées visibles, un appel à l'action clair.

Maintenant que nous savons concevoir un flyer, nous allons, dans la section suivante, apprendre à \terme{créer des visuels pour les réseaux sociaux}.

% ============================================================
\section{Création de visuels pour les réseaux sociaux}
% ============================================================

\subsection{L'importance des réseaux sociaux}

Les \terme{réseaux sociaux} sont devenus des canaux de communication incontournables pour les marques, les entreprises, les associations, les créateurs. Ils permettent de :

\begin{itemize}
    \item \terme{Toucher} un large public.
    \item \terme{Créer} une communauté.
    \item \terme{Promouvoir} un produit, un service, un événement.
    \item \terme{Construire} une identité de marque.
    \item \terme{Engager} les abonnés.
\end{itemize}

Les visuels jouent un rôle \terme{essentiel} sur les réseaux sociaux : ils attirent l'attention, transmettent un message, créent une émotion, incitent à l'action.

\subsection{Les spécificités des réseaux sociaux}

\subsubsection{Les formats}

Chaque plateforme a ses formats spécifiques :

\begin{table}[h!]
\centering
\begin{tabular}{@{}p{3cm}p{4cm}p{6cm}@{}}
\toprule
\textbf{Plateforme} & \textbf{Type} & \textbf{Dimensions (px)} \\
\midrule
Instagram & Post carré & 1080 × 1080 \\
Instagram & Post portrait & 1080 × 1350 \\
Instagram & Story & 1080 × 1920 \\
Instagram & Reel & 1080 × 1920 \\
Facebook & Post & 1200 × 630 \\
Facebook & Story & 1080 × 1920 \\
Facebook & Couverture & 851 × 315 \\
Twitter / X & Post & 1200 × 675 \\
LinkedIn & Post & 1200 × 627 \\
LinkedIn & Couverture & 1584 × 396 \\
YouTube & Miniature & 1280 × 720 \\
Pinterest & Épingle & 1000 × 1500 \\
TikTok & Vidéo & 1080 × 1920 \\
\bottomrule
\end{tabular}
\caption{Dimensions des visuels pour les réseaux sociaux.}
\end{table}

\subsubsection{Les contraintes}

\begin{itemize}
    \item \terme{Taille} : les visuels sont souvent vus en petit (smartphone).
    \item \terme{Lisibilité} : le texte doit être lisible sur petit écran.
    \item \terme{Rapidité} : les visuels sont vus en quelques secondes.
    \item \terme{Concurrence} : l'attention est difficile à capter.
    \item \terme{Zone de sécurité} : certains éléments peuvent être masqués par l'interface.
\end{itemize}

\subsection{La conception des visuels}

\subsubsection{Le message}

Sur les réseaux sociaux, le message doit être :

\begin{itemize}
    \item \terme{Court} : 3 à 7 mots pour le texte principal.
    \item \terme{Percutant} : fort, mémorable.
    \item \terme{Lisible} : grand, contrasté.
    \item \terme{Visuel} : privilégier l'image au texte.
\end{itemize}

\subsubsection{La composition}

La composition doit être :

\begin{itemize}
    \item \terme{Simple} : pas de surcharge.
    \item \terme{Équilibrée} : éléments bien répartis.
    \item \terme{Hiérarchisée} : message principal, secondaire, tertiaire.
    \item \terme{Adaptée} : au format (carré, portrait, story).
\end{itemize}

\subsubsection{La couleur}

Les couleurs doivent être :

\begin{itemize}
    \item \terme{Cohérentes} : avec la charte graphique.
    \item \terme{Contrastées} : pour attirer l'attention.
    \item \terme{Expressives} : pour créer une émotion.
    \item \terme{Lisibles} : le texte doit contraster avec le fond.
\end{itemize}

\subsubsection{La typographie}

La typographie doit être :

\begin{itemize}
    \item \terme{Lisible} : taille suffisante.
    \item \terme{Limitée} : 1 à 2 polices.
    \item \terme{Cohérente} : avec la charte graphique.
    \item \terme{Adaptée} : au format.
\end{itemize}

\subsection{Les types de visuels}

\subsubsection{Le visuel de post}

Le \terme{visuel de post} est l'image publiée sur le fil d'actualité. Il doit attirer l'attention, transmettre le message, inciter à l'engagement.

\subsubsection{La story}

La \terme{story} est un format vertical, éphémère (24h). Elle doit être :

\begin{itemize}
    \item \terme{Verticale} : adaptée au format 9:16.
    \item \terme{Dynamique} : avec du mouvement, des animations.
    \item \terme{Interactive} : avec des sondages, des questions.
    \item \terme{Cohérente} : avec l'identité de marque.
\end{itemize}

\subsubsection{La couverture}

La \terme{couverture} est l'image qui orne le haut d'un profil (Facebook, LinkedIn, YouTube). Elle doit :

\begin{itemize}
    \item \terme{Respecter} les dimensions.
    \item \terme{Prévoir} une zone de sécurité.
    \item \terme{Cohérente} avec l'identité de marque.
\end{itemize}

\subsubsection{La bannière}

La \terme{bannière} est une image promotionnelle. Elle peut être utilisée pour :

\begin{itemize}
    \item \terme{Promouvoir} un événement.
    \item \terme{Annoncer} une offre.
    \item \terme{Mettre} en avant un produit.
\end{itemize}

\subsection{Les bonnes pratiques}

\begin{itemize}
    \item \terme{Respecter} les dimensions de chaque plateforme.
    \item \terme{Prévoir} une zone de sécurité.
    \item \terme{Assurer} la lisibilité sur petit écran.
    \item \terme{Privilégier} l'image au texte.
    \item \terme{Utiliser} des couleurs contrastées.
    \item \terme{Limiter} le nombre de polices.
    \item \terme{Créer} une série cohérente (même style, mêmes couleurs).
    \item \terme{Adapter} le format (carré, portrait, story).
    \item \terme{Tester} sur mobile avant publication.
\end{itemize}

\subsection*{Conclusion de la section}

Les \terme{visuels pour les réseaux sociaux} sont essentiels pour la communication des marques. Chaque plateforme a ses formats spécifiques. Les contraintes sont la taille, la lisibilité, la rapidité, la concurrence, la zone de sécurité. La conception doit privilégier un message court, une composition simple, des couleurs cohérentes, une typographie lisible. Il est essentiel de respecter les dimensions, de prévoir une zone de sécurité, d'assurer la lisibilité, de créer une série cohérente, de tester sur mobile.

Maintenant que nous savons créer des visuels pour les réseaux sociaux, nous allons, dans la section suivante, apprendre à \terme{réaliser une affiche événementielle}.

% ============================================================
\section{Réalisation d'une affiche événementielle}
% ============================================================

\subsection{Qu'est-ce qu'une affiche événementielle ?}

\begin{quote}
\textbf{Une affiche événementielle est une affiche qui annonce un événement : concert, festival, exposition, conférence, spectacle, compétition sportive. Elle doit informer, attirer, donner envie de venir.}
\end{quote}

L'affiche événementielle est un support \terme{essentiel} pour la communication d'un événement. Elle doit :

\begin{itemize}
    \item \terme{Attirer} l'attention dans un environnement chargé.
    \item \terme{Communiquer} l'ambiance de l'événement.
    \item \terme{Informer} sur les éléments essentiels (nom, date, lieu).
    \item \terme{Donner} envie de venir.
\end{itemize}

\subsection{Les éléments d'une affiche événementielle}

\subsubsection{Le nom de l'événement}

Le \terme{nom} est l'élément principal. Il doit être :

\begin{itemize}
    \item \terme{Grand} : visible de loin.
    \item \terme{Lisible} : police adaptée, contraste.
    \item \terme{Expressif} : typographie en cohérence avec l'ambiance.
\end{itemize}

\subsubsection{La date et l'heure}

La \terme{date} et l'\terme{heure} doivent être visibles, mais secondaires :

\begin{itemize}
    \item \terme{Taille} : moyenne.
    \item \terme{Position} : souvent en bas, parfois en haut.
    \item \terme{Lisibilité} : contraste suffisant.
\end{itemize}

\subsubsection{Le lieu}

Le \terme{lieu} doit être précis : nom du lieu, adresse, ville. Il est souvent placé avec la date.

\subsubsection{Le visuel}

Le \terme{visuel} transmet l'ambiance de l'événement :

\begin{itemize}
    \item \terme{Photographie} : artiste, lieu, ambiance.
    \item \terme{Illustration} : créative, expressive.
    \item \terme{Graphique} : formes, couleurs, typographie.
\end{itemize}

\subsubsection{Les artistes/intervenants}

Pour un concert, un festival, un spectacle, les \terme{artistes} sont des éléments essentiels. Ils sont souvent listés sur l'affiche.

\subsubsection{Le logo}

Le \terme{logo} de l'organisateur ou du sponsor identifie l'événement.

\subsubsection{Les informations pratiques}

Les \terme{informations pratiques} : tarifs, réservation, site web, réseaux sociaux.

\subsubsection{Les mentions légales}

Les \terme{mentions légales} : organisateur, licences, etc.

\subsection{Les styles d'affiche événementielle}

\subsubsection{Le style photographique}

Utilise une \terme{photographie} forte : artiste, lieu, ambiance. Convient aux concerts, festivals, spectacles.

\subsubsection{Le style illustratif}

Utilise une \terme{illustration} créative : dessin, collage, graphisme. Convient aux festivals, expositions, événements culturels.

\subsubsection{Le style typographique}

Utilise la \terme{typographie} comme élément principal : compositions de mots, formes géométriques. Convient aux événements modernes, artistiques.

\subsubsection{Le style minimaliste}

Utilise peu d'éléments : fond uni, typographie épurée. Convient aux événements élégants, haut de gamme.

\subsection{La conception d'une affiche événementielle}

\subsubsection{Étape 1 : Comprendre l'événement}

\begin{itemize}
    \item \terme{Type} : concert, festival, exposition, conférence.
    \item \terme{Public} : âge, milieu, habitudes.
    \item \terme{Ambiance} : festive, élégante, intimiste, dynamique.
\end{itemize}

\subsubsection{Étape 2 : Choisir le style}

Choisir le style adapté à l'événement : photographique, illustratif, typographique, minimaliste.

\subsubsection{Étape 3 : Définir la hiérarchie}

Définir l'ordre d'importance :

\begin{itemize}
    \item \terme{Nom de l'événement} (1).
    \item \terme{Visuel} (2).
    \item \terme{Artistes} (3).
    \item \terme{Date et lieu} (4).
    \item \terme{Informations pratiques} (5).
    \item \terme{Logo} (6).
    \item \terme{Mentions légales} (7).
\end{itemize}

\subsubsection{Étape 4 : Composer}

Composer l'affiche :

\begin{itemize}
    \item \terme{Placer} le visuel.
    \item \terme{Placer} le nom de l'événement.
    \item \terme{Placer} les artistes.
    \item \terme{Placer} la date et le lieu.
    \item \terme{Placer} les informations pratiques.
    \item \terme{Placer} le logo.
    \item \terme{Placer} les mentions légales.
\end{itemize}

\subsubsection{Étape 5 : Appliquer les couleurs et la typographie}

\begin{itemize}
    \item \terme{Couleurs} : cohérentes avec l'ambiance.
    \item \terme{Typographie} : cohérente avec l'ambiance, lisible.
\end{itemize}

\subsubsection{Étape 6 : Vérifier}

\begin{itemize}
    \item \terme{Lisibilité} : de loin, de près.
    \item \terme{Cohérence} : visuel, couleurs, typographie.
    \item \terme{Impact} : attire l'attention.
    \item \terme{Complétude} : toutes les informations sont présentes.
\end{itemize}

\subsection{Les bonnes pratiques}

\begin{itemize}
    \item \terme{Un visuel fort} et cohérent.
    \item \terme{Un nom d'événement} lisible et expressif.
    \item \terme{Une hiérarchie} claire.
    \item \terme{Des informations} complètes et lisibles.
    \item \terme{Une ambiance} cohérente.
    \item \terme{Une typographie} limitée.
    \item \terme{Des couleurs} cohérentes.
    \item \terme{Vérifier} la lisibilité de loin.
    \item \terme{Tester} sur différents supports.
\end{itemize}

\subsection*{Conclusion de la section}

L'\terme{affiche événementielle} annonce un événement : concert, festival, exposition, conférence. Elle doit informer, attirer, donner envie de venir. Les éléments sont le nom, la date, le lieu, le visuel, les artistes, le logo, les informations pratiques, les mentions légales. Les styles sont photographique, illustratif, typographique, minimaliste. La conception suit plusieurs étapes : comprendre l'événement, choisir le style, définir la hiérarchie, composer, appliquer les couleurs et la typographie, vérifier. Il est essentiel d'avoir un visuel fort, une hiérarchie claire, des informations complètes, une ambiance cohérente.

Maintenant que nous savons réaliser une affiche événementielle, nous allons, dans la section suivante, apprendre à \terme{créer une bannière publicitaire}.

% ============================================================
\section{Création d'une bannière publicitaire}
% ============================================================

\subsection{Qu'est-ce qu'une bannière publicitaire ?}

\begin{quote}
\textbf{Une bannière publicitaire est un visuel promotionnel destiné à être affiché sur le web, dans des espaces publicitaires, ou sur les réseaux sociaux. Elle doit attirer l'attention, transmettre un message, inciter au clic.}
\end{quote}

Les bannières sont utilisées pour :

\begin{itemize}
    \item \terme{Promouvoir} un produit, un service, un événement.
    \item \terme{Annoncer} une offre, une promotion.
    \item \terme{Développer} la notoriété d'une marque.
    \item \terme{Générer} du trafic vers un site web.
    \item \terme{Convertir} des visiteurs en clients.
\end{itemize}

\subsection{Les formats de bannière}

\subsubsection{Les formats standards}

\begin{table}[h!]
\centering
\begin{tabular}{@{}p{4cm}p{4cm}p{5cm}@{}}
\toprule
\textbf{Format} & \textbf{Dimensions (px)} & \textbf{Usage} \\
\midrule
Leaderboard & 728 × 90 & Haut de page \\
Medium Rectangle & 300 × 250 & Contenu, sidebar \\
Large Rectangle & 336 × 280 & Contenu, sidebar \\
Skyscraper & 120 × 600 & Sidebar \\
Wide Skyscraper & 160 × 600 & Sidebar \\
Half Page & 300 × 600 & Sidebar \\
Billboard & 970 × 250 & Haut de page \\
Mobile Leaderboard & 320 × 50 & Mobile \\
Mobile Banner & 320 × 100 & Mobile \\
\bottomrule
\end{tabular}
\caption{Formats standards de bannières publicitaires.}
\end{table}

\subsubsection{Les formats pour les réseaux sociaux}

\begin{itemize}
    \item \terme{Facebook} : 1200 × 628.
    \item \terme{Instagram} : 1080 × 1080 (carré), 1080 × 1350 (portrait).
    \item \terme{Twitter / X} : 1200 × 675.
    \item \terme{LinkedIn} : 1200 × 627.
\end{itemize}

\subsection{Les éléments d'une bannière}

\subsubsection{Le visuel}

Le \terme{visuel} attire l'attention : produit, personne, illustration, graphique.

\subsubsection{Le titre}

Le \terme{titre} transmet le message principal : court, percutant, lisible.

\subsubsection{Le sous-titre}

Le \terme{sous-titre} complète le titre : argument, avantage, offre.

\subsubsection{L'appel à l'action}

L'\terme{appel à l'action} (CTA) indique l'action à effectuer : « Acheter », « Découvrir », « S'inscrire », etc. Il doit être visible, distinctif.

\subsubsection{Le logo}

Le \terme{logo} identifie la marque. Il doit être visible mais discret.

\subsubsection{Les mentions légales}

Les \terme{mentions légales} : conditions, mentions obligatoires.

\subsection{La conception d'une bannière}

\subsubsection{Étape 1 : Définir l'objectif}

Quel est l'objectif de la bannière ? (notoriété, conversion, information)

\subsubsection{Étape 2 : Définir le message}

Quel est le message principal ? Court, percutant, lisible.

\subsubsection{Étape 3 : Choisir le visuel}

Choisir un visuel attractif, adapté au message et au format.

\subsubsection{Étape 4 : Composer}

\begin{itemize}
    \item \terme{Placer} le visuel.
    \item \terme{Placer} le titre.
    \item \terme{Placer} le sous-titre.
    \item \terme{Placer} le CTA.
    \item \terme{Placer} le logo.
    \item \terme{Placer} les mentions légales.
\end{itemize}

\subsubsection{Étape 5 : Appliquer les couleurs et la typographie}

\begin{itemize}
    \item \terme{Couleurs} : cohérentes avec la marque, contrastées.
    \item \terme{Typographie} : lisible, limitée.
\end{itemize}

\subsubsection{Étape 6 : Vérifier}

\begin{itemize}
    \item \terme{Lisibilité} : à petite taille.
    \item \terme{Impact} : attire l'attention.
    \item \terme{Clarté} : message compris en 3 secondes.
    \item \terme{Cohérence} : avec la marque.
\end{itemize}

\subsection{Les bonnes pratiques}

\begin{itemize}
    \item \terme{Message} court et clair.
    \item \terme{Visuel} attractif.
    \item \terme{Texte} lisible à petite taille.
    \item \terme{CTA} visible.
    \item \terme{Logo} présent.
    \item \terme{Couleurs} cohérentes et contrastées.
    \item \terme{Typographie} limitée.
    \item \terme{Poids} optimisé (pour le web).
    \item \terme{Tester} sur différents formats.
    \item \terme{Décliner} en plusieurs tailles.
\end{itemize}

\subsection*{Conclusion de la section}

La \terme{bannière publicitaire} est un visuel promotionnel destiné au web ou aux réseaux sociaux. Elle doit attirer l'attention, transmettre un message, inciter au clic. Les formats sont variés (leaderboard, rectangle, skyscraper, etc.). Les éléments sont le visuel, le titre, le sous-titre, le CTA, le logo, les mentions légales. La conception suit plusieurs étapes : définir l'objectif, définir le message, choisir le visuel, composer, appliquer les couleurs et la typographie, vérifier. Il est essentiel d'avoir un message court, un visuel attractif, un CTA visible, des couleurs cohérentes.

Maintenant que nous savons créer une bannière publicitaire, nous allons, dans la section suivante, apprendre à \terme{intégrer des photographies et des textes}.

% ============================================================
\section{Intégration des photographies et des textes}
% ============================================================

\subsection{L'importance de l'intégration}

L'\terme{intégration} des photographies et des textes est cruciale dans un support publicitaire. Une photographie mal intégrée (mauvaise résolution, mauvais cadrage, mauvaise couleur) et un texte mal intégré (illisible, mal placé, incohérent) ruinent le visuel. Une intégration soignée donne un résultat professionnel.

\subsection{L'intégration des photographies}

\subsubsection{La résolution}

Les photographies doivent avoir une \terme{résolution suffisante} :

\begin{itemize}
    \item \terme{Impression} : 300 ppi à la taille finale.
    \item \terme{Web} : 72 ppi (ou 144 ppi pour Retina).
\end{itemize}

\subsubsection{Le cadrage}

Le \terme{cadrage} de la photographie doit être adapté à la composition :

\begin{itemize}
    \item \terme{Recadrer} pour se concentrer sur le sujet.
    \item \terme{Recadrer} pour respecter la règle des tiers.
    \item \terme{Recadrer} pour respecter le format du support.
\end{itemize}

\subsubsection{La couleur}

La \terme{couleur} de la photographie doit être harmonisée avec la composition :

\begin{itemize}
    \item \terme{Ajuster} la balance des couleurs.
    \item \terme{Ajuster} la saturation.
    \item \terme{Ajuster} le contraste.
\end{itemize}

\subsubsection{Le détourage}

Si nécessaire, \terme{détourer} la photographie pour l'intégrer dans la composition.

\subsubsection{L'application d'effets}

On peut appliquer des effets à la photographie :

\begin{itemize}
    \item \terme{Flou} : pour créer une profondeur.
    \item \terme{Ombre} : pour détacher l'image.
    \item \terme{Vignettage} : pour attirer le regard au centre.
    \item \terme{Texture} : pour créer un effet.
\end{itemize}

\subsection{L'intégration des textes}

\subsubsection{La lisibilité}

Les textes doivent être \terme{lisibles} :

\begin{itemize}
    \item \terme{Contraste} : le texte doit contraster avec le fond.
    \item \terme{Taille} : suffisamment grand.
    \item \terme{Police} : adaptée au support et au public.
\end{itemize}

\subsubsection{La hiérarchie}

Les textes doivent être \terme{hiérarchisés} :

\begin{itemize}
    \item \terme{Titre} : le plus grand, le plus visible.
    \item \terme{Sous-titre} : secondaire.
    \item \terme{Texte courant} : plus petit.
    \item \terme{Légendes} : encore plus petit.
    \item \terme{Mentions} : le plus petit.
\end{itemize}

\subsubsection{La position}

Les textes doivent être \terme{bien positionnés} :

\begin{itemize}
    \item \terme{Alignement} : avec les autres éléments.
    \item \terme{Proximité} : avec les éléments liés.
    \item \terme{Respiration} : espaces autour du texte.
\end{itemize}

\subsubsection{Les effets}

On peut appliquer des effets au texte :

\begin{itemize}
    \item \terme{Ombre portée} : pour détacher le texte.
    \item \terme{Contour} : pour améliorer la lisibilité.
    \item \terme{Dégradé} : pour ajouter de la couleur.
    \item \terme{Texture} : pour ajouter de la matière.
\end{itemize}

\subsection{L'intégration du texte et de l'image}

\subsubsection{Le principe}

Le texte et l'image doivent \terme{s'intégrer} :

\begin{itemize}
    \item \terme{Le texte} ne doit pas être « posé » sur l'image, mais intégré.
    \item \terme{L'image} ne doit pas être « décorée » par le texte, mais accompagnée.
    \item \terme{Les deux} doivent former un tout cohérent.
\end{itemize}

\subsubsection{Les techniques}

\begin{itemize}
    \item \terme{Masque} : masquer une partie du texte avec l'image (ou inversement).
    \item \terme{Mode de fusion} : utiliser un mode de fusion pour intégrer le texte.
    \item \terme{Ombre} : ajouter une ombre pour détacher le texte.
    \item \terme{Contour} : ajouter un contour pour améliorer la lisibilité.
    \item \terme{Dégradé} : ajouter un dégradé pour intégrer le texte.
\end{itemize}

\subsection{Les bonnes pratiques}

\begin{itemize}
    \item \terme{Choisir} des photographies de qualité.
    \item \terme{Recadrer} les photographies avec soin.
    \item \terme{Harmoniser} les couleurs.
    \item \terme{Détourer} si nécessaire.
    \item \terme{Assurer} la lisibilité du texte.
    \item \terme{Hiérarchiser} les textes.
    \item \terme{Positionner} les textes avec soin.
    \item \terme{Appliquer} les effets avec parcimonie.
    \item \terme{Intégrer} texte et image.
    \item \terme{Vérifier} le rendu final.
\end{itemize}

\subsection*{Conclusion de la section}

L'\terme{intégration des photographies et des textes} est cruciale dans un support publicitaire. On choisit des photographies de qualité, on les recadre, on harmonise les couleurs, on détoure si nécessaire. On assure la lisibilité du texte, on hiérarchise, on positionne, on applique des effets avec parcimonie. On intègre texte et image pour former un tout cohérent. Il est essentiel de vérifier le rendu final.

Maintenant que nous savons intégrer des photographies et des textes, nous allons, dans la section suivante, apprendre à \terme{organiser la hiérarchie visuelle}.

% ============================================================
\section{Organisation de la hiérarchie visuelle}
% ============================================================

\subsection{L'importance de la hiérarchie visuelle}

La \terme{hiérarchie visuelle} organise les éléments d'un support publicitaire selon leur importance. Elle permet :

\begin{itemize}
    \item \terme{Guider} le regard du spectateur.
    \item \terme{Transmettre} le message dans le bon ordre.
    \item \terme{Faciliter} la compréhension.
    \item \terme{Mettre} en valeur les éléments clés.
\end{itemize}

Sans hiérarchie, le spectateur ne sait pas où regarder, le message est confus, l'efficacité est nulle.

\subsection{Les niveaux de hiérarchie}

\subsubsection{Niveau 1 : le message principal}

Le \terme{message principal} est l'élément le plus important : titre, visuel, accroche. Il doit être vu en premier.

\subsubsection{Niveau 2 : les messages secondaires}

Les \terme{messages secondaires} complètent le message principal : sous-titre, visuel secondaire, description courte.

\subsubsection{Niveau 3 : les informations pratiques}

Les \terme{informations pratiques} : date, lieu, prix, coordonnées, appel à l'action.

\subsubsection{Niveau 4 : les mentions}

Les \terme{mentions} : mentions légales, crédits, copyright.

\subsection{Les moyens de créer la hiérarchie}

\subsubsection{La taille}

Plus un élément est grand, plus il est important.

\subsubsection{La graisse}

Plus un élément est gras, plus il est important.

\subsubsection{La couleur}

Les couleurs vives et contrastées attirent plus que les couleurs neutres.

\subsubsection{La position}

Les éléments en haut, au centre, à gauche sont vus en premier (dans les cultures occidentales).

\subsubsection{L'isolement}

Un élément isolé par du vide attire plus.

\subsubsection{Le contraste}

Les éléments contrastés attirent plus.

\subsubsection{Les effets}

Les éléments avec effets (ombre, contour, dégradé) attirent plus.

\subsection{Construire une hiérarchie}

\subsubsection{Étape 1 : Identifier les niveaux}

Déterminer les niveaux de la composition : message principal, messages secondaires, informations pratiques, mentions.

\subsubsection{Étape 2 : Attribuer les moyens}

Attribuer à chaque niveau les moyens adaptés : taille, graisse, couleur, position, effets.

\subsubsection{Étape 3 : Vérifier la cohérence}

Vérifier que la hiérarchie est cohérente sur l'ensemble de la composition.

\subsubsection{Étape 4 : Tester}

Tester la composition : le message principal est-il vu en premier ? Les informations sont-elles lisibles ?

\subsection{Les erreurs de hiérarchie}

\subsubsection{Pas de hiérarchie}

Tous les éléments ont la même importance. Le spectateur ne sait pas où regarder.

\subsubsection{Hiérarchie inversée}

Le message secondaire est plus visible que le message principal.

\subsubsection{Trop de niveaux}

Trop de tailles, de graisses, de couleurs différentes. Le spectateur est perdu.

\subsubsection{Hiérarchie incohérente}

Les niveaux ne sont pas respectés.

\subsection{Les bonnes pratiques}

\begin{itemize}
    \item \terme{Un seul message principal}.
    \item \terme{Peu de niveaux} (3 à 4).
    \item \terme{Moyens adaptés} à chaque niveau.
    \item \terme{Cohérence} sur l'ensemble de la composition.
    \item \terme{Vérifier} la lisibilité.
    \item \terme{Tester} avec un œil neuf.
\end{itemize}

\subsection*{Conclusion de la section}

La \terme{hiérarchie visuelle} organise les éléments d'un support publicitaire selon leur importance. Elle comporte généralement trois ou quatre niveaux : message principal, messages secondaires, informations pratiques, mentions. Elle se crée par la taille, la graisse, la couleur, la position, l'isolement, le contraste, les effets. Il est essentiel d'avoir un seul message principal, peu de niveaux, des moyens adaptés, une cohérence, et de vérifier le résultat.

Maintenant que nous savons organiser la hiérarchie visuelle, nous allons, dans la section suivante, apprendre à \terme{décliner un visuel en plusieurs formats}.

% ============================================================
\section{Déclinaison d'un visuel en plusieurs formats}
% ============================================================

\subsection{Pourquoi décliner un visuel ?}

La \terme{déclinaison} consiste à adapter un visuel à plusieurs formats. Elle permet :

\begin{itemize}
    \item \terme{Communiquer} sur plusieurs supports (imprimé, web, réseaux sociaux).
    \item \terme{Toucher} plusieurs publics.
    \item \terme{Assurer} la cohérence de la communication.
    \item \terme{Optimiser} l'impact de la campagne.
    \item \terme{Gagner} du temps (créer une fois, décliner plusieurs fois).
\end{itemize}

La déclinaison est une pratique courante en agence : une campagne est déclinée sur affiche, flyer, bannière, réseaux sociaux, etc.

\subsection{Les supports à décliner}

\subsubsection{Les supports imprimés}

\begin{itemize}
    \item \terme{Affiche} A3, A2, A1, A0.
    \item \terme{Flyer} A5, A6, DL.
    \item \terme{Dépliant} A4 3 volets.
    \item \terme{Carte de visite} 85 × 55 mm.
    \item \terme{Papier à en-tête} A4.
\end{itemize}

\subsubsection{Les supports numériques}

\begin{itemize}
    \item \terme{Site web} : bannière, hero.
    \item \terme{Réseaux sociaux} : post, story, couverture.
    \item \terme{Newsletter} : bandeau, visuel.
    \item \terme{Publicité en ligne} : leaderboard, rectangle.
\end{itemize}

\subsection{La méthode de déclinaison}

\subsubsection{Étape 1 : Créer le visuel principal}

Créer le visuel dans le format le plus grand (par exemple, affiche A3). Ce visuel sert de référence.

\subsubsection{Étape 2 : Identifier les éléments à conserver}

Identifier les éléments essentiels : titre, visuel, logo, couleurs, typographies. Ces éléments doivent être conservés dans toutes les déclinaisons.

\subsubsection{Étape 3 : Identifier les éléments à adapter}

Identifier les éléments à adapter : position, taille, format. Ces éléments doivent être adaptés à chaque format.

\subsubsection{Étape 4 : Créer les déclinaisons}

Créer les déclinaisons en adaptant les éléments :

\begin{itemize}
    \item \terme{Changer} les dimensions du document.
    \item \terme{Repositionner} les éléments.
    \item \terme{Redimensionner} les éléments.
    \item \terme{Adapter} les marges.
\end{itemize}

\subsubsection{Étape 5 : Vérifier la cohérence}

Vérifier que toutes les déclinaisons sont cohérentes : mêmes couleurs, mêmes typographies, même style.

\subsubsection{Étape 6 : Exporter}

Exporter chaque déclinaison dans le bon format et la bonne résolution.

\subsection{Les bonnes pratiques}

\begin{itemize}
    \item \terme{Créer} le visuel dans le format le plus grand.
    \item \terme{Conserver} les éléments essentiels (titre, logo, couleurs).
    \item \terme{Adapter} les éléments à chaque format.
    \item \terme{Respecter} les proportions.
    \item \terme{Adapter} les marges.
    \item \terme{Vérifier} la cohérence.
    \item \terme{Vérifier} la lisibilité dans chaque format.
    \item \terme{Exporter} dans le bon format.
    \item \terme{Organiser} les fichiers (par format, par version).
\end{itemize}

\subsection{Les erreurs courantes}

\subsubsection{Ne pas adapter le visuel}

Utiliser le même visuel dans tous les formats, sans adaptation, donne un résultat médiocre.

\subsubsection{Perdre la cohérence}

Les déclinaisons n'ont plus la même identité (couleurs, typographies, style).

\subsubsection{Oublier des éléments}

Certaines déclinaisons oublient des éléments essentiels (logo, mentions).

\subsubsection{Ne pas vérifier}

Ne pas vérifier la lisibilité, la cohérence, la qualité de chaque déclinaison.

\subsection*{Conclusion de la section}

La \terme{déclinaison d'un visuel} consiste à adapter un visuel à plusieurs formats. Elle permet de communiquer sur plusieurs supports, de toucher plusieurs publics, d'assurer la cohérence. Les supports sont imprimés (affiche, flyer, dépliant) ou numériques (site web, réseaux sociaux, bannière). La méthode consiste à créer le visuel principal, identifier les éléments à conserver et à adapter, créer les déclinaisons, vérifier la cohérence, exporter. Il est essentiel de conserver les éléments essentiels, d'adapter les éléments, de vérifier la cohérence.

Maintenant que nous savons décliner un visuel, nous allons, dans la section suivante, apprendre à \terme{contrôler la qualité}.

% ============================================================
\section{Contrôle de la qualité}
% ============================================================

\subsection{L'importance du contrôle qualité}

Le \terme{contrôle qualité} est une étape essentielle avant la livraison d'un support publicitaire. Il permet de :

\begin{itemize}
    \item \terme{Détecter} les erreurs.
    \item \terme{Corriger} les problèmes.
    \item \terme{Assurer} la qualité du résultat.
    \item \terme{Éviter} les mauvaises surprises (impression, affichage, web).
    \item \terme{Préserver} la réputation du designer.
\end{itemize}

Un support publicitaire mal contrôlé peut ruiner une campagne, une marque, une réputation.

\subsection{Les points à contrôler}

\subsubsection{La lisibilité}

\begin{itemize}
    \item \terme{Texte} : lisible à la taille finale.
    \item \terme{Contraste} : suffisant entre texte et fond.
    \item \terme{Taille} : adaptée au support.
    \item \terme{Police} : adaptée au public.
\end{itemize}

\subsubsection{La hiérarchie}

\begin{itemize}
    \item \terme{Message principal} : visible en premier.
    \item \terme{Messages secondaires} : visibles en second.
    \item \terme{Informations pratiques} : lisibles.
    \item \terme{Mentions} : présentes.
\end{itemize}

\subsubsection{Les couleurs}

\begin{itemize}
    \item \terme{Cohérence} : avec la charte graphique.
    \item \terme{Harmonie} : couleurs qui s'accordent.
    \item \terme{Contraste} : suffisant.
    \item \terme{Profil} : sRGB (web), CMJN (impression).
\end{itemize}

\subsubsection{Les images}

\begin{itemize}
    \item \terme{Résolution} : suffisante.
    \item \terme{Netteté} : correcte.
    \item \terme{Couleur} : harmonisée.
    \item \terme{Cadrage} : adapté.
    \item \terme{Droits} : libres de droits.
\end{itemize}

\subsubsection{La typographie}

\begin{itemize}
    \item \terme{Polices} : limitées (1 à 2).
    \item \terme{Lisibilité} : bonne.
    \item \terme{Alignement} : correct.
    \item \terme{Orthographe} : sans fautes.
    \item \terme{Ponctuation} : correcte.
\end{itemize}

\subsubsection{Les éléments techniques}

\begin{itemize}
    \item \terme{Marges} : suffisantes.
    \item \terme{Fonds perdus} : présents (impression).
    \item \terme{Traits de coupe} : présents (impression).
    \item \terme{Polices} : incorporées (PDF).
    \item \terme{Liens} : corrects (web).
\end{itemize}

\subsubsection{Le contenu}

\begin{itemize}
    \item \terme{Message} : clair, cohérent.
    \item \terme{Informations} : complètes, exactes.
    \item \terme{Mentions légales} : présentes.
    \item \terme{Appel à l'action} : clair.
\end{itemize}

\subsection{La méthode de contrôle}

\subsubsection{Étape 1 : Vérifier le contenu}

Vérifier le contenu : textes, informations, mentions.

\subsubsection{Étape 2 : Vérifier la lisibilité}

Zoomer, dézoomer, imprimer, tester sur mobile.

\subsubsection{Étape 3 : Vérifier la hiérarchie}

Vérifier que le message principal est vu en premier.

\subsubsection{Étape 4 : Vérifier les couleurs}

Vérifier la cohérence, l'harmonie, le profil.

\subsubsection{Étape 5 : Vérifier les images}

Vérifier la résolution, la netteté, la couleur, le cadrage.

\subsubsection{Étape 6 : Vérifier la typographie}

Vérifier les polices, la lisibilité, l'alignement, l'orthographe.

\subsubsection{Étape 7 : Vérifier les éléments techniques}

Vérifier les marges, les fonds perdus, les traits de coupe, les polices incorporées.

\subsubsection{Étape 8 : Vérifier sur différents supports}

Vérifier sur écran, sur papier (impression), sur mobile.

\subsection{Les bonnes pratiques}

\begin{itemize}
    \item \terme{Vérifier} chaque point systématiquement.
    \item \terme{Utiliser} une checklist.
    \item \terme{Zoomer} pour voir les détails.
    \item \terme{Imprimer} pour vérifier le rendu.
    \item \terme{Tester} sur mobile.
    \item \terme{Faire} relire par quelqu'un d'autre.
    \item \terme{Corriger} les erreurs.
    \item \terme{Re-vérifier} après correction.
\end{itemize}

\subsection*{Conclusion de la section}

Le \terme{contrôle qualité} est essentiel avant la livraison d'un support publicitaire. Il permet de détecter les erreurs, de corriger les problèmes, d'assurer la qualité. Les points à contrôler sont la lisibilité, la hiérarchie, les couleurs, les images, la typographie, les éléments techniques, le contenu. La méthode consiste à vérifier systématiquement chaque point, à zoomer, à imprimer, à tester sur mobile, à faire relire. Il est essentiel de vérifier, de corriger, de re-vérifier.

Maintenant que nous savons contrôler la qualité, nous allons, dans la dernière section de ce chapitre, apprendre à \terme{exporter les fichiers finaux}.

% ============================================================
\section{Exportation des fichiers finaux}
% ============================================================

\subsection{L'importance de l'exportation}

L'\terme{exportation} est l'étape finale de la création d'un support publicitaire. Elle consiste à préparer le fichier pour son usage final (impression, web, réseaux sociaux). Une exportation mal faite peut ruiner tout le travail : fichier illisible, couleurs fausses, résolution insuffisante, poids excessif.

\subsection{Les formats d'exportation}

\subsubsection{Pour l'impression}

\begin{itemize}
    \item \terme{PDF/X-1a} : standard pour l'impression, polices incorporées, CMJN.
    \item \terme{PDF/X-4} : standard moderne, transparence gérée.
    \item \terme{TIFF} : image haute qualité, CMJN, 300 ppi.
    \item \terme{EPS} : format vectoriel, pour certains imprimeurs.
\end{itemize}

\subsubsection{Pour le web}

\begin{itemize}
    \item \terme{JPEG} : photographies, compression avec perte.
    \item \terme{PNG} : images avec transparence, graphiques.
    \item \terme{WebP} : format moderne, plus léger.
    \item \terme{SVG} : vectoriel, pour les logos et icônes.
\end{itemize}

\subsubsection{Pour les réseaux sociaux}

\begin{itemize}
    \item \terme{JPEG} : post, story, couverture.
    \item \terme{PNG} : images avec transparence.
    \item \terme{MP4} : vidéos, reels.
\end{itemize}

\subsubsection{Pour l'archivage}

\begin{itemize}
    \item \terme{PSD} : fichier source Photoshop.
    \item \terme{TIFF} : image haute qualité.
    \item \terme{PDF} : document.
\end{itemize}

\subsection{Les paramètres d'exportation}

\subsubsection{Pour l'impression}

\begin{itemize}
    \item \terme{Résolution} : 300 ppi.
    \item \terme{Profil colorimétrique} : CMJN (selon l'imprimeur).
    \item \terme{Fonds perdus} : 3 à 5 mm.
    \item \terme{Traits de coupe} : présents.
    \item \terme{Polices} : incorporées.
    \item \terme{Images} : incorporées ou liées.
\end{itemize}

\subsubsection{Pour le web}

\begin{itemize}
    \item \terme{Résolution} : 72 ppi (ou 144 ppi pour Retina).
    \item \terme{Profil colorimétrique} : sRGB.
    \item \terme{Dimensions} : adaptées au support.
    \item \terme{Compression} : adaptée (JPEG qualité 60-80, PNG optimisé).
    \item \terme{Poids} : optimisé (moins de 200 ko pour les bannières).
\end{itemize}

\subsubsection{Pour les réseaux sociaux}

\begin{itemize}
    \item \terme{Résolution} : 72 ppi.
    \item \terme{Profil colorimétrique} : sRGB.
    \item \terme{Dimensions} : spécifiques à chaque plateforme.
    \item \terme{Compression} : adaptée.
    \item \terme{Poids} : optimisé.
\end{itemize}

\subsection{La méthode d'exportation}

\subsubsection{Étape 1 : Enregistrer le fichier source}

Toujours conserver le fichier source (PSD) avant d'exporter.

\subsubsection{Étape 2 : Choisir le format}

Choisir le format adapté à l'usage : PDF, TIFF (impression), JPEG, PNG (web), JPEG (réseaux sociaux).

\subsubsection{Étape 3 : Configurer les paramètres}

Configurer les paramètres : résolution, profil colorimétrique, dimensions, compression.

\subsubsection{Étape 4 : Vérifier le rendu}

Vérifier le rendu du fichier exporté.

\subsubsection{Étape 5 : Nommer le fichier}

Nommer le fichier de manière claire et descriptive.

\subsubsection{Étape 6 : Organiser}

Organiser les fichiers exportés dans des dossiers clairs.

\subsection{Les bonnes pratiques}

\begin{itemize}
    \item \terme{Conserver} le fichier source (PSD).
    \item \terme{Choisir} le bon format selon l'usage.
    \item \terme{Configurer} les bons paramètres.
    \item \terme{Vérifier} le rendu.
    \item \terme{Nommer} les fichiers clairement.
    \item \terme{Organiser} les fichiers.
    \item \terme{Optimiser} le poids (web).
    \item \terme{Fournir} les fichiers au bon format au client.
\end{itemize}

\subsection{Les erreurs courantes}

\subsubsection{Mauvais format}

Exporter en JPEG pour l'impression (perte de qualité), en PNG pour les photos (poids excessif).

\subsubsection{Mauvaise résolution}

Exporter en 72 ppi pour l'impression (pixellisation), en 300 ppi pour le web (poids excessif).

\subsubsection{Mauvais profil}

Exporter en RVB pour l'impression (couleurs fausses), en CMJN pour le web (couleurs fausses).

\subsubsection{Fichier trop lourd}

Exporter sans optimiser le poids (web), le fichier met du temps à charger.

\subsubsection{Fichier mal nommé}

Exporter sans nommage clair, difficile à retrouver.

\subsection*{Conclusion de la section}

L'\terme{exportation} est l'étape finale de la création d'un support publicitaire. Elle consiste à préparer le fichier pour son usage final. Les formats sont : PDF, TIFF (impression) ; JPEG, PNG, WebP (web) ; JPEG, PNG, MP4 (réseaux sociaux) ; PSD, TIFF, PDF (archivage). Les paramètres sont : résolution, profil colorimétrique, dimensions, compression. Il est essentiel de conserver le fichier source, de choisir le bon format, de configurer les bons paramètres, de vérifier le rendu, de nommer clairement, d'organiser.

% ============================================================
\section*{Conclusion du chapitre}
\addcontentsline{toc}{section}{Conclusion du chapitre}
% ============================================================

Nous voici au terme du chapitre 13. Nous avons appris à créer des supports publicitaires avec Photoshop :

\begin{itemize}
    \item \terme{Méthodologie de création d'une affiche} : brief, recherche, conception, création, présentation, finalisation, production, évaluation ; contraintes (format, résolution, colorimétrie, marges).
    \item \terme{Conception d'un flyer commercial} : formats, recto/verso, structure (recto attractif, verso informatif).
    \item \terme{Création de visuels pour les réseaux sociaux} : dimensions par plateforme, contraintes, conception, types (post, story, couverture, bannière).
    \item \terme{Réalisation d'une affiche événementielle} : éléments (nom, date, lieu, visuel, artistes, logo, informations), styles (photographique, illustratif, typographique, minimaliste).
    \item \terme{Création d'une bannière publicitaire} : formats standards, éléments (visuel, titre, sous-titre, CTA, logo, mentions).
    \item \terme{Intégration des photographies et des textes} : résolution, cadrage, couleur, détourage, lisibilité, hiérarchie, position, effets.
    \item \terme{Organisation de la hiérarchie visuelle} : niveaux (message principal, secondaires, informations, mentions), moyens (taille, graisse, couleur, position, isolement, contraste, effets).
    \item \terme{Déclinaison d'un visuel} : supports (imprimés, numériques), méthode (créer le visuel principal, adapter les éléments, vérifier la cohérence).
    \item \terme{Contrôle de la qualité} : lisibilité, hiérarchie, couleurs, images, typographie, éléments techniques, contenu.
    \item \terme{Exportation des fichiers finaux} : formats (PDF, TIFF, JPEG, PNG, WebP), paramètres (résolution, profil, dimensions, compression), bonnes pratiques.
\end{itemize}

La création de supports publicitaires est l'aboutissement de la maîtrise de Photoshop. Elle mobilise toutes les compétences acquises : création de documents, organisation des calques, sélection, détourage, retouche, photomontage, texte, effets. Elle nécessite aussi une compréhension du contexte : brief, public, objectifs, contraintes. Un designer qui maîtrise Photoshop mais ne comprend pas le contexte produira des visuels vides ; un designer qui comprend le contexte et maîtrise Photoshop produira des visuels efficaces.

Ce chapitre t'a donné les bases de la création de supports publicitaires. Mais la maîtrise vient avec la pratique. N'hésite pas à t'entraîner sur des projets variés : affiches, flyers, visuels réseaux sociaux, bannières. Chaque projet est un nouveau défi, et chaque défi est une occasion de progresser.

Ce chapitre clôt la deuxième partie de ce cours, consacrée à la \terme{maîtrise d'Adobe Photoshop}. Tu disposes désormais des compétences nécessaires pour créer, retoucher, détourer, photomonter, composer, exporter des visuels professionnels. Ces compétences sont essentielles, mais elles ne suffisent pas : il faut aussi maîtriser les images vectorielles, qui sont utilisées pour les logos, les illustrations, les mises en page. C'est l'objet de la partie suivante, consacrée à Adobe Illustrator.

Dans le chapitre suivant, nous aborderons \terme{Adobe Illustrator}. Tu découvriras ce logiciel, ses domaines d'utilisation, ses différences avec Photoshop, son interface, ses outils. Tu apprendras à créer des documents, à dessiner des formes vectorielles, à utiliser la plume, à manipuler les formes, à gérer les couleurs et les dégradés, à créer du texte, à concevoir des logos et des identités visuelles. Ces compétences te permettront de compléter ta maîtrise du design graphique.

% ============================================================
\section*{Projet de la Partie II}
\addcontentsline{toc}{section}{Projet de la Partie II}
% ============================================================

\subsection*{Concevoir une affiche publicitaire complète}

\subsubsection*{Objectif}

Mettre en pratique l'ensemble des compétences acquises dans la Partie II : création de documents, organisation des calques, sélection, détourage, retouche photographique, photomontage, texte, effets graphiques, exportation.

\subsubsection*{Sujet}

Concevoir une affiche publicitaire complète pour un produit, un service, un événement ou une marque fictive. Le sujet est libre, mais il doit être réaliste et cohérent.

\subsubsection*{Contraintes}

\begin{itemize}
    \item \terme{Format} : A3 (29,7 × 42 cm), 300 ppi, CMJN.
    \item \terme{Contenu} : un titre (message principal), un sous-titre, un visuel (photographie retouchée et/ou photomontage), un texte (informations), un logo, un appel à l'action, des mentions légales.
    \item \terme{Techniques} : au moins une retouche photographique, un détourage, un photomontage, une typographie travaillée, des effets graphiques.
    \item \terme{Organisation} : calques nommés, groupés, colorés ; travail non destructif.
\end{itemize}

\subsubsection*{Étapes}

\begin{enumerate}
    \item \terme{Brief} : définir le produit, le public, le message, l'objectif.
    \item \terme{Recherche} : références, moodboard, esquisses.
    \item \terme{Création} : document, images, détourage, retouche, photomontage, texte, effets.
    \item \terme{Présentation} : préparer une présentation du projet (démarche, choix).
    \item \terme{Finalisation} : contrôle qualité, exportation.
\end{enumerate}

\subsubsection*{Auto-évaluation}

\begin{table}[h!]
\centering
\begin{tabular}{@{}lccccc@{}}
\toprule
\textbf{Critère} & \textbf{1} & \textbf{2} & \textbf{3} & \textbf{4} & \textbf{5} \\
\midrule
Brief compris et respecté & \(\square\) & \(\square\) & \(\square\) & \(\square\) & \(\square\) \\
Recherche pertinente & \(\square\) & \(\square\) & \(\square\) & \(\square\) & \(\square\) \\
Détourage soigné & \(\square\) & \(\square\) & \(\square\) & \(\square\) & \(\square\) \\
Retouche naturelle & \(\square\) & \(\square\) & \(\square\) & \(\square\) & \(\square\) \\
Photomontage cohérent & \(\square\) & \(\square\) & \(\square\) & \(\square\) & \(\square\) \\
Typographie maîtrisée & \(\square\) & \(\square\) & \(\square\) & \(\square\) & \(\square\) \\
Effets maîtrisés & \(\square\) & \(\square\) & \(\square\) & \(\square\) & \(\square\) \\
Hiérarchie claire & \(\square\) & \(\square\) & \(\square\) & \(\square\) & \(\square\) \\
Organisation des calques & \(\square\) & \(\square\) & \(\square\) & \(\square\) & \(\square\) \\
Exportation correcte & \(\square\) & \(\square\) & \(\square\) & \(\square\) & \(\square\) \\
Qualité globale & \(\square\) & \(\square\) & \(\square\) & \(\square\) & \(\square\) \\
\bottomrule
\end{tabular}
\caption{Grille d'auto-évaluation du projet de la Partie II.}
\end{table}

\subsubsection*{Livrables}

\begin{itemize}
    \item \terme{Fichier PSD} : fichier source complet, avec calques organisés.
    \item \terme{Fichier PDF} : version imprimable, CMJN, 300 ppi.
    \item \terme{Fichier JPEG} : version web, sRGB, 72 ppi.
    \item \terme{Présentation} : document de présentation (démarche, choix, visuels).
\end{itemize}

\vspace{2em}

\begin{center}
\textit{Fin du chapitre 13 --- Création de supports publicitaires avec Photoshop.}

\vspace{1em}

\textit{Fin de la Partie II --- Maîtrise d'Adobe Photoshop.}

\vspace{1em}

\textbf{Prochain chapitre : Chapitre 14 --- Introduction à Adobe Illustrator.}
\end{center}